\begin{table}[htbp]
\centering
\caption{Post-hoc compactness audit of \textit{baimu fang} formed by the block-level DRL agent. 
For each township and each evaluation seed we recover the final farmland configuration 
by replaying the agent's block selection sequence, identify all connected farmland 
components with area $\geq$6.67\,ha (the raw \textit{baimu fang} count), and compute 
the isoperimetric quotient $\mathrm{IPQ} = 4\pi A / P^2$ on the dissolved farmland polygon 
of each component (IPQ$=$1 for a circle; lower means more elongated/fragmented shape). 
The ``Init.\ IPQ'' columns report the same compactness counts on the \emph{initial} land-use 
configuration before any optimization, providing a fairness baseline: the Queen-contiguity 
connectivity definition combined with real cadastral parcel shapes typically yields 
low-IPQ components even prior to consolidation, so the post-DRL counts at each 
$\tau$ should be read against the corresponding initial value rather than against an ideal of $\tau{=}1$. 
DRL values are means over 5 seeds.}
\label{tab:compactness}
\resizebox{\textwidth}{!}{%
\begin{tabular}{lrrrrrrrrr}
\toprule
 & \multicolumn{4}{c}{Initial state} & \multicolumn{5}{c}{After DRL optimization (mean over 5 seeds)} \\
\cmidrule(lr){2-5}\cmidrule(lr){6-10}
Township & Init.\ raw & Init.\ IPQ$\geq$0.10 & Init.\ IPQ$\geq$0.20 & Init.\ IPQ$\geq$0.30 & Final raw & Final IPQ$\geq$0.10 & Final IPQ$\geq$0.20 & Final IPQ$\geq$0.30 & $\Delta$ raw \\
\midrule
A & 9 & 3 & 0 & 0 & 9.0 & 0.6 & 0.0 & 0.0 & +0.0 \\
B & 6 & 1 & 0 & 0 & 5.0 & 1.6 & 0.0 & 0.0 & -1.0 \\
C & 6 & 0 & 0 & 0 & 11.0 & 1.0 & 0.0 & 0.0 & +5.0 \\
\bottomrule
\end{tabular}%
}
\end{table}
